
% --------------------------------------------------------------- CONFIGURATIONS

%ifdef TWOSIDE
	%\documentclass[a4paper,12pt,final,twoside,openright]{book}
%elif ONESIDE
	\documentclass[a4paper,12pt,final,oneside]{book}
%endif

\usepackage{rapport}


% -------------------------------------------------------------- META: CONSTANTS

\newcommand{\reporttitle}{ERP}
\newcommand{\enseignants}{Youssef~\textsc{Amghar}\\ Anne~\textsc{Legait}\\ Pierre-Alain~\textsc{Millet}\\ Mohamed~\textsc{Ouhalima}}
\newcommand{\reportauthor}{Guillaume~\textsc{Abadie}\\ Thierry~\textsc{Cantenot}\\ Marina~\textsc{Julien}\\ Ahmed~\textsc{Kachkach}\\ Martin~\textsc{Wetterwald}\\ Valentina~\textsc{Zantedeschi}}
\newcommand{\hexanome}{SpecIFic}
\newcommand{\reportsubject}{Livrable de projet}
\newcommand{\stagetopic}{Blablabla test}
\newcommand{\dateperiod}{du 17 décembre 2013 au 11 mars 2014}
\newcommand{\HRule}{\rule{\linewidth}{0.5mm}}
\setlength{\parskip}{1ex} % Espace entre les paragraphes

\hypersetup{
	pdftitle={\reporttitle},%
		pdfauthor={\reportauthor},%
		pdfsubject={\reportsubject},%
		pdfkeywords={INSA Lyon}
}

\title{\reporttitle}
\author{\reportauthor}
%\setcounter{tocdepth}{4}


% ------------------------------------------------------------------------- FILE

\begin{document}


    % ------------------------------------------------------------------- HEADER

	\renewcommand{\chaptername}{} %\renewcommand{\thechapter}{}
	\renewcommand{\contentsname}{Sommaire}

	\pagestyle{empty}
	\pagenumbering{Roman}


    % ------------------------------------------------------------ HEADER: TITLE

	% Inspiré de http://en.wikibooks.org/wiki/LaTeX/Title_Creation
\begin{center}
	\begin{minipage}[t]{0.48\textwidth}
	  \begin{flushleft}
	    \includegraphics [width=40mm]{images/logo_INSA.png} \\[0.5cm]
			INSA Lyon\\
			20, avenue Albert Einstein\\
			69621 Villeurbanne Cedex
	  \end{flushleft}
	\end{minipage}
	\begin{minipage}[t]{0.48\textwidth}
	  \begin{flushright}
	    %\includegraphics [width=60mm]{images/logo_Passau.jpg} \\[0.5cm]
	    %Universität Passau\\
		%Innstraße, 3\\
		%	D-94032 Passau
	  \end{flushright}
	\end{minipage} \\[2cm]

	\textsc{\Large \reportsubject}\\[0.3cm]
	\HRule \\[0.4cm]
	{\Huge \bfseries \reporttitle}\\[0.3cm]
	{\LARGE \bfseries «~\stagetopic~»}\\[0.3cm]
	{\Large \dateperiod}\\[0.4cm]
	\HRule \\[1cm]

	\includegraphics [scale=0.35]{images/i7.jpg} \\[0.7cm]
	\begin{minipage}[t]{0.4\textwidth}
	  \begin{flushleft} \large
	    \emph{Hexanôme \textbf{\hexanome}~:}\\
	    \small \reportauthor
	  \end{flushleft}
	\end{minipage}
	\begin{minipage}[t]{0.5\textwidth}
	  \begin{flushright} \large
	    \emph{Enseignants~:} \\
	    \enseignants
	  \end{flushright}
	\end{minipage}

	\vfill
	\footnotesize{Année scolaire 2013-2014}
\end{center}


	%ifdef TWOSIDE
		%\cleardoublepage
	%endif

	%\include{title2}

	%ifdef TWOSIDE
		%  \newpage
		%	\null
		%	\vfill
	%endif


    % --------------------------------------------------- HEADER: CONFIGURATIONS

	\sloppy          % Justification moins stricte : des mots ne dépasseront pas des paragraphes

    \frontmatter
		\pagestyle{empty}
		\tableofcontents
		\addtocontents{toc}{\protect\thispagestyle{empty}}

	\mainmatter
	\pagestyle{headings}

	\renewcommand{\chaptermark}[1]{\markboth{\MakeUppercase{\chaptername\ \thechapter.\ #1}}{}}
	\renewcommand{\sectionmark}[1]{\markright{\thesection{} #1}}


    % ------------------------------------------------------------------ CONTENT

	\chapter*{Introduction}
\addcontentsline{toc}{chapter}{Introduction}
\chaptermark{Introduction}
Ceci est un test de pipeautage.


	\chapter*{Conclusion}
\addcontentsline{toc}{chapter}{Conclusion}
\chaptermark{Conclusion}
C'était un test de pipeautage.


	%\renewcommand{\chaptermark}[1]{\markboth{\MakeUppercase{#1}}{}}
	%\renewcommand{\sectionmark}[1]{\markright{#1}}

	%\addcontentsline{toc}{part}{Annexes}
	%\part*{Annexes}
	%\appendix
	%\include{implementationExercices}


    % ------------------------------------------------------------------- FOOTER
\end{document}
